\documentclass[11pt,a4paper]{article}
\usepackage[utf8]{inputenc}
\usepackage[spanish,es-nodecimaldot]{babel}
\usepackage[margin=2cm]{geometry}
\usepackage{xltabular}
\usepackage{booktabs}
\usepackage{hyperref}
\usepackage{microtype}
\usepackage{url}

\hypersetup{
    colorlinks=true,
    linkcolor=blue,
    filecolor=magenta,      
    urlcolor=blue,
    citecolor=blue,
    pdftitle={Actividad Formativa FORM-TL-01: Taller de tipos de requisitos},
    pdfauthor={Luis Alberto Paca Pilataza}
}

\title{\textbf{Universidad Técnica Estatal de Quevedo}\\[0.3em]
\large Facultad de Ciencias de la Computación, Carrera de Software\\[0.5em]
\Large Ingeniería de Requerimientos (ISR-401)\\[0.8em]
\LARGE \textbf{Actividad Formativa FORM-TL-01}\\
\Large Taller de tipos de requisitos}

\author{\textbf{Estudiante:} Luis Alberto Paca Pilataza \\
\textbf{Curso:} 4.\textsuperscript{o} Nivel Software, Paralelo A}
\date{Octubre 2026}

\begin{document}

\maketitle

\section{Clasificación de Enunciados (Tabla 1)}

A continuación se presenta la clasificación de los 20 enunciados del caso de estudio «Policlínico Ribera del Quevedo», evaluados en sus tres dimensiones fundamentales: tipo de requisito (funcional, calidad según ISO/IEC 25010:2023 o restricción), alcance (producto o proceso) y nivel (sistema o software), acompañados de una justificación clara y directa en una sola línea.

\begin{small}
\begin{xltabular}{\textwidth}{|c|X|p{3.4cm}|c|c|X|}
\hline
\textbf{Cód.} & \textbf{Enunciado (Resumido)} & \textbf{Dim. 1: Tipo} & \textbf{Dim. 2} & \textbf{Dim. 3} & \textbf{Justificación} \\ \hline
\endfirsthead

\hline
\textbf{Cód.} & \textbf{Enunciado (Resumido)} & \textbf{Dim. 1: Tipo} & \textbf{Dim. 2} & \textbf{Dim. 3} & \textbf{Justificación} \\ \hline
\endhead

\hline
\multicolumn{6}{|r|}{\textit{Continúa en la siguiente página\dots}} \\ \hline
\endfoot

\hline
\endlastfoot

E01 & El sistema debe permitir al paciente reservar citas médicas eligiendo especialidad, médico y horario. & Funcional & Producto & Software & Es un servicio directo de reserva que el programa debe ofrecer al paciente en pantalla. \\ \hline

E02 & El sistema debe enviar un mensaje de texto de recordatorio al paciente 24 horas antes de su cita. & Funcional & Producto & Sistema & Requiere que el programa se conecte con una red celular externa para entregar el mensaje. \\ \hline

E03 & La aplicación debe ser rápida. & Calidad (Eficiencia de desempeño, Comportamiento temporal) & Producto & Software & Expresa una expectativa sobre el tiempo de respuesta del programa sin definir una tarea específica. \\ \hline

E04 & El sistema debe permitir al recepcionista registrar en ventanilla la llegada presencial del paciente. & Funcional & Producto & Sistema & Apoya una tarea de atención presencial donde interactúan la persona en ventanilla y el flujo de la clínica. \\ \hline

E05 & La plataforma se construirá con lenguaje Java 21 y base de datos PostgreSQL. & Restricción (Tecnológica, Arquitectura) & Producto & Software & Obliga a usar un lenguaje y una base de datos específicos para construir el programa. \\ \hline

E06 & El sistema debe permitir al médico consultar la lista diaria de citas y marcar las consultas atendidas. & Funcional & Producto & Software & Define opciones directas de consulta y cambio de estado dentro de la agenda del médico. \\ \hline

E07 & Los datos de salud de los pacientes deben estar seguros. & Calidad (Seguridad, Confidencialidad e Integridad) & Producto & Software & Pide proteger la información médica para que no sea vista ni modificada por personas no autorizadas. \\ \hline

E08 & La aplicación móvil debe funcionar en teléfonos con sistema operativo Android 10 o superior. & Restricción (Entorno operativo) & Producto & Sistema & Fija los modelos de celular y versiones de sistema operativo donde debe funcionar la aplicación. \\ \hline

E09 & El sistema debe registrar como inasistencia toda cita cancelada con menos de cuatro horas de anticipación. & Funcional (Regla de negocio) & Producto & Software & Establece una regla interna de atención que el programa debe aplicar de forma automática. \\ \hline

E10 & La plataforma debe ser fácil de usar para los pacientes adultos mayores. & Calidad (Capacidad de interacción, Aprendizaje y Accesibilidad) & Producto & Sistema & Describe la facilidad y comodidad de uso orientada a las capacidades de las personas de la tercera edad. \\ \hline

E11 & El proyecto completo debe entregarse en un plazo máximo de seis meses. & Restricción (Temporal, Calendario) & Proceso & Sistema & Fija una fecha límite para el calendario de trabajo del equipo, sin relación con el código en sí. \\ \hline

E12 & El sistema debe cerrar la sesión del usuario de forma automática tras 10 minutos de inactividad. & Funcional (Control de sesión) & Producto & Software & Indica la acción que tomará el programa para proteger la cuenta cuando el usuario deje de usarla. \\ \hline

E13 & El sistema debe permitir leer la cédula en pantalla táctil e imprimir un ticket de turno en el quiosco. & Funcional & Producto & Sistema & Detalla un servicio de recepción física que combina la pantalla táctil y la impresora del quiosco. \\ \hline

E14 & La plataforma debe estar siempre disponible. & Calidad (Fiabilidad, Disponibilidad) & Producto & Sistema & Requiere que el servicio web y móvil esté accesible y funcionando durante el horario de atención. \\ \hline

E15 & El sistema debe permitir al director médico generar reportes mensuales de inasistencias por especialidad. & Funcional & Producto & Software & Es una función que permite al directivo consultar y descargar resúmenes de citas no asistidas. \\ \hline

E16 & El equipo debe validar cada incremento de software con el director médico cada dos semanas. & Restricción (Metodológica) & Proceso & Sistema & Define la forma de trabajo del equipo mediante revisiones quincenales con el directivo. \\ \hline

E17 & El manejo de información de pacientes debe cumplir la Ley Orgánica de Protección de Datos Personales. & Restricción (Legal, Regulatoria) & Producto & Sistema & Exige respetar las leyes ecuatorianas de protección de datos al manejar información de los pacientes. \\ \hline

E18 & La plataforma debe integrarse sin problemas con el sistema de facturación actual. & Calidad (Compatibilidad, Interoperabilidad) & Producto & Sistema & Solicita que el nuevo programa logre comunicarse y compartir datos con el sistema de cobros actual. \\ \hline

E19 & El sistema debe restablecer el servicio en menos de 30 minutos sin perder citas tras fallar el servidor. & Calidad (Fiabilidad, Recuperabilidad) & Producto & Sistema & Fija el tiempo máximo de recuperación ante caídas del servidor y asegura que no se pierdan citas guardadas. \\ \hline

E20 & Todo cambio en el código fuente debe ser revisado por otro integrante del equipo antes de integrarse. & Restricción (Procedimental, Control de calidad) & Proceso & Software & Establece una regla interna de control para que ningún programador suba cambios sin la revisión de un compañero. \\ \hline

\end{xltabular}
\end{small}

\section{Cuantificación de Requisitos de Calidad}

Siguiendo los lineamientos de la norma ISO/IEC 25010:2023 \cite{iso25010} y las recomendaciones de ingeniería de requisitos \cite{pohl2025, iso29148}, se identificaron cinco enunciados ambiguos del caso de estudio (E03, E07, E10, E14 y E18). A continuación se transforman en requisitos de calidad verificables, definiendo para cada uno su característica de calidad, métrica, umbral exacto fundamentado en los datos del policlínico y método de verificación.

\begin{enumerate}
    \item \textbf{Enunciado original (E03):} «La aplicación debe ser \textit{rápida}».
    \begin{itemize}
        \item \textbf{Característica ISO/IEC 25010:2023:} Eficiencia de desempeño, Comportamiento temporal.
        \item \textbf{Requisito cuantificado (RNF-01):} En horarios de máxima demanda (lunes a las 07:00 con hasta 300 usuarios conectados de manera simultánea), la plataforma deberá cargar y mostrar la disponibilidad de citas médicas en un tiempo no mayor a 2,0 segundos para el 95\% de las solicitudes realizadas desde la web o la aplicación móvil.
        \item \textbf{Métrica y umbral:} Tiempo de respuesta en la consulta de turnos $\le 2,0$ segundos (percentil 95) bajo carga de 300 usuarios concurrentes.
        \item \textbf{Método de verificación:} Pruebas de rendimiento y estrés simulando 300 conexiones simultáneas contra la base de datos de citas.
    \end{itemize}

    \item \textbf{Enunciado original (E07):} «Los datos de salud de los pacientes deben estar \textit{seguros}».
    \begin{itemize}
        \item \textbf{Característica ISO/IEC 25010:2023:} Seguridad, Confidencialidad e Integridad.
        \item \textbf{Requisito cuantificado (RNF-02):} El 100\% de los datos médicos y personales de los pacientes almacenados y transmitidos deberán protegerse mediante protocolos de comunicación cifrada (TLS 1.3) y cifrado en base de datos (AES-256). Solo el médico con cita asignada y el propio paciente tendrán acceso a visualizar el detalle de la consulta médica, en conformidad con la Ley Orgánica de Protección de Datos Personales del Ecuador \cite{asamblea2021}.
        \item \textbf{Métrica y umbral:} Porcentaje de registros médicos cifrados = 100\%; intentos no autorizados permitidos = 0\%.
        \item \textbf{Método de verificación:} Auditoría técnica de seguridad sobre la base de datos e inspección de políticas de acceso por roles.
    \end{itemize}

    \item \textbf{Enunciado original (E10):} «La plataforma debe ser \textit{fácil de usar} para los pacientes adultos mayores».
    \begin{itemize}
        \item \textbf{Característica ISO/IEC 25010:2023:} Capacidad de interacción, Facilidad de aprendizaje y Accesibilidad.
        \item \textbf{Requisito cuantificado (RNF-03):} Al menos el 85\% de una muestra de pacientes mayores de 60 años deberá lograr reservar una cita médica desde un teléfono con Android 10 en un tiempo menor o igual a 4 minutos en su primer intento, sin requerir asistencia de otra persona.
        \item \textbf{Métrica y umbral:} Tasa de finalización exitosa de la reserva $\ge 85\%$; tiempo por intento $\le 4$ minutos.
        \item \textbf{Método de verificación:} Prueba práctica de usabilidad en la sala de espera con un grupo piloto de 20 pacientes mayores de 60 años.
    \end{itemize}

    \item \textbf{Enunciado original (E14):} «La plataforma debe estar \textit{siempre disponible}».
    \begin{itemize}
        \item \textbf{Característica ISO/IEC 25010:2023:} Fiabilidad, Disponibilidad.
        \item \textbf{Requisito cuantificado (RNF-04):} La plataforma deberá mantener una disponibilidad mensual mínima del 99,8\% dentro del horario de atención del policlínico (lunes a sábado de 07:00 a 19:00), tolerando un tiempo máximo acumulado de caída no programada de 36 minutos al mes.
        \item \textbf{Métrica y umbral:} Porcentaje de operatividad mensual en horario de consulta $\ge 99,8\%$ (caída máxima no programada $\le 36$ minutos al mes).
        \item \textbf{Método de verificación:} Reportes mensuales automáticos emitidos por herramientas de monitoreo continuo del servidor web.
    \end{itemize}

    \item \textbf{Enunciado original (E18):} «La plataforma debe integrarse \textit{sin problemas} con el sistema de facturación actual».
    \begin{itemize}
        \item \textbf{Característica ISO/IEC 25010:2023:} Compatibilidad, Interoperabilidad.
        \item \textbf{Requisito cuantificado (RNF-05):} La plataforma deberá comunicarse con el servicio web del sistema de facturación existente para registrar cada cobro generado, confirmando la transacción en menos de 2,0 segundos y alcanzando una tasa de éxito mínima del 99,9\% en el envío de datos válidos, reintentando automáticamente en caso de fallas de red.
        \item \textbf{Métrica y umbral:} Tiempo de respuesta de integración $\le 2,0$ segundos; tasa de éxito en el registro de cobros $\ge 99,9\%$.
        \item \textbf{Método de verificación:} Pruebas de integración automatizadas entre el módulo de reservas y el servicio web de facturación.
    \end{itemize}
\end{enumerate}

\section{Discusión de Discrepancias con Compañero}

Como parte de la actividad, se realizó una sesión de comparación de resultados con un compañero de curso. Al contrastar nuestras tablas iniciales, encontramos dos discrepancias de criterio que analizamos y resolvimos con base en la teoría estudiada:

\subsection*{Discrepancia 1: Enunciado E02 (Recordatorio por mensaje de texto 24 horas antes)}
\begin{itemize}
    \item \textbf{Postura inicial del compañero:} Clasificó el enunciado como requisito de \textbf{Software}, argumentando que el temporizador de 24 horas y la redacción del mensaje se programan en el código fuente.
    \item \textbf{Mi postura:} Sostuve que pertenecía a la dimensión de \textbf{Sistema}.
    \item \textbf{Resolución acordada:} Al revisar la teoría de Klaus Pohl \cite{pohl2025} y el estándar IEEE~29148 \cite{iso29148}, observamos que cuando una función depende indispensablemente de infraestructura de comunicaciones externa (red celular y pasarela SMS) y del teléfono del paciente para completarse en el mundo real, supera el límite del software puro y se ubica en el nivel de sistema. Por ende, acordamos mantener la clasificación como requisito de \textbf{Sistema}.
\end{itemize}

\subsection*{Discrepancia 2: Enunciado E10 (Fácil de usar para pacientes adultos mayores)}
\begin{itemize}
    \item \textbf{Postura inicial del compañero:} Consideró inicialmente que se trataba de una \textbf{Restricción de Proceso}, bajo la idea de que obligaba a los desarrolladores a cumplir pautas de diseño durante la fase de maquetación.
    \item \textbf{Mi postura:} Defendí que correspondía a un \textbf{Requisito de Calidad del Producto}.
    \item \textbf{Resolución consensuada:} Consultamos la norma ISO/IEC 25010:2023 \cite{iso25010}, la cual cataloga la «Capacidad de interacción» (Interaction capability) como una característica intrínseca de calidad del producto software. La facilidad de uso evalúa cómo el usuario final experimenta la plataforma terminada y no una regla sobre cómo gestionar el proyecto. Con base en esto, convenimos que su clasificación correcta es \textbf{Requisito de Calidad del Producto}.
\end{itemize}

\section{Declaración de Uso de Inteligencia Artificial}

En cumplimiento de la Norma 8 de la actividad formativa y del Código de Ética de la ACM \cite{acm2018}, se detalla de manera transparente el uso de herramientas de inteligencia artificial durante la realización de este taller:

\begin{itemize}
    \item \textbf{Herramienta utilizada:} Gemini.
    \item \textbf{Propósito exclusivo de apoyo:} Se utilizó Gemini para corregir errores ortográficos, revisar signos de puntuación (como comas y puntos), mejorar la estructura gramatical y organizar el formato del código en \LaTeX.
    \item \textbf{Forma de verificación aplicada por el estudiante:} Todo el análisis técnico, la identificación de los tipos de requisitos y la clasificación de los 20 enunciados fueron revisados y validados personalmente por el autor con base en las normas ISO/IEC 25010:2023 e ISO/IEC/IEEE 29148. Asimismo, los valores numéricos y umbrales de los requisitos de calidad fueron calculados directamente a partir de los datos del caso de estudio, comprobando finalmente la compilación limpia del documento en \texttt{pdfLaTeX}.
\end{itemize}

\begin{thebibliography}{5}

\bibitem{pohl2025}
K.~Pohl, \emph{Requirements Engineering: Fundamentals, Principles, and Techniques}, 2nd~ed.\hskip 1em plus
  0.5em minus 0.4em\relax Berlin, Heidelberg, Germany: Springer, 2025. ISBN 978-3-662-69204-2.

\bibitem{iso29148}
ISO/IEC/IEEE, \emph{ISO/IEC/IEEE 29148:2018. Systems and software engineering, Life cycle processes, Requirements engineering}.\hskip 1em plus
  0.5em minus 0.4em\relax International Standard, 2018.

\bibitem{iso25010}
ISO/IEC, \emph{ISO/IEC 25010:2023. Systems and software engineering, Systems and software quality requirements and evaluation (SQuaRE), Product quality model}.\hskip 1em plus
  0.5em minus 0.4em\relax International Standard, 2023.

\bibitem{asamblea2021}
Asamblea Nacional del Ecuador, \emph{Ley Orgánica de Protección de Datos Personales}.\hskip 1em plus
  0.5em minus 0.4em\relax Registro Oficial, Quinto Suplemento No. 459, 26 de mayo de 2021, 2021.

\bibitem{acm2018}
Association for Computing Machinery, \emph{ACM Code of Ethics and Professional Conduct}, 2018. [En línea]. Disponible en: \url{https://www.acm.org/code-of-ethics}. Consultado el 27 de octubre de 2026.

\end{thebibliography}

\end{document}
